First suppose both $a,b>1$. Coprimality forces, in some order, $a=2^u,b=3^v$ with $u,v\ge1$. Then $n$ is odd. Squares modulo 8 show that $u=1,2$ are impossible and that for $u\ge3$, $v$ must be even. Write $v=2w$. Factoring gives
\[(n-3^w)(n+3^w)=2^u.\]
Both factors are even, and their difference is $2\cdot3^w$, so the smaller must be 2. If the larger is $2^s$, then $2^{s-1}=3^w+1$. For even $w>0$, the right side is 2 modulo 8, impossible; for odd $w$, it is 4 modulo 8, forcing $s-1=2,w=1$. This gives $n=5$.

Now suppose $a=1$, so $(n-1)(n+1)=2^u3^v$, allowing zero exponents. If $n$ is even, the two odd factors are coprime powers of 3, forcing $n-1=1$ and $n=2$. If $n$ is odd, put $r=\dfrac{n-1}{2}$ and $s=\dfrac{n+1}{2}$. They are coprime consecutive integers with only prime divisors 2 and 3. If $r=1$, then $n=3$. Otherwise they are a positive power of 2 and a positive power of 3 differing by one.

The only such pairs are $\{2,3\},\{3,4\},\{8,9\}$. Indeed, $2^k=3^\ell+1$ gives $(k,\ell)=(2,1)$ by the modulo-8 argument above. For $3^\ell=2^k+1$, $k=1$ gives $(1,1)$; otherwise reduction modulo 3 forces odd $k\ge3$, and reduction modulo 8 makes $\ell$ even. The two powers of 2 in $(3^{\dfrac{\ell}{2}}-1)(3^{\dfrac{\ell}{2}}+1)=2^k$ differ by 2, so they are 2 and 4, giving $(k,\ell)=(3,2)$. Therefore the odd values here are $n=5,7,17$.

All $n\in\{2,3,5,7,17\}$ occur, using $(a,b)=(1,3),(1,8),(1,24),(1,48),(1,288)$. Their product is $\boxed{3570}$.
